\guidelinesubsubsection{Rationale}

Transparency about LLM involvement is a prerequisite for informed assessment of a study's scope, limitations, and potential biases.
Without explicit disclosure, readers cannot evaluate how the LLM's characteristics may have influenced the research process or its outcomes.

\guidelinesubsubsection{Recommendations}

When conducting any kind of empirical study involving LLMs, researchers \must clearly declare that an LLM was used (see \scope for what we consider relevant research support).
This \should be done in a suitable section of the \paper, for example, in the introduction or research methods section; \citeauthor{cheng2025writing}~\cite{cheng2025writing} argue specifically for the methods section, since acknowledgments are easily missed at the end of a paper.
% GROUNDING cheng2025writing: "the most helpful and transparent means for describing the use of these tools is within the methods section of a manuscript"
The ACM Policy on Authorship requires authors to describe in the methods section any use of AI in the research itself~\cite{ACM2026}.
% GROUNDING ACM2026: "the specific use(s) of AI tools must be described in detail in the methods section of the Work"

Beyond generic declarations, researchers \should report the exact purpose of using an LLM in a study, the tasks it was used to automate, and the expected benefits in the \paper.
A sufficient declaration specifies not only \emph{that} an LLM was used, but also \emph{which} LLM (name and version), \emph{how} it was used (e.g., as an annotator, code generator, or judge), and \emph{where} in the research process it was employed (e.g., data collection, analysis, or synthesis).

When the LLM is central to the study (e.g., as the main tool being evaluated or as a core component of the research method), the declaration \should be prominent and detailed, appearing in the methodology section with cross-references to the specific guidelines that apply (e.g., Sections \modelversion, \design, and \traces).
When the LLM's role is more tangential (e.g., used for a single preprocessing step), a brief but explicit statement in the methodology section is sufficient.
When a study assigns multiple distinct roles to LLMs (e.g., one model generates evaluation data while another scores outputs), each role \should be declared separately.
In each case, the disclosure \must be specific enough for readers to assess how the LLM's involvement may affect the study's validity and reproducibility.

\guidelinesubsubsection{Examples}

The \emph{ACM Policy on Authorship} requires reporting AI-generated artifacts such as code, datasets, and figures where they underlie a study's conclusions~\cite{ACM2026}.
% GROUNDING ACM2026: "the creation of artifacts that are directly relevant to the conclusions of the research, such as code, datasets, and charts or figures that rely on the AI tools"
Reporting in the methods section keeps this disclosure visible under double-blind review, where end-of-paper acknowledgments are often removed.
An example of an LLM disclosure beyond writing support can be found in a recent paper by \citeauthor{DBLP:conf/re/LubosFTGMEL24}~\cite{DBLP:conf/re/LubosFTGMEL24}, in which they write in the methodology section:
% GROUNDING DBLP:conf/re/LubosFTGMEL24: "We conducted an LLM-based evaluation of requirements utilizing the Llama 2 language model with 70 billion parameters, fine-tuned to complete chat responses."

\begin{quote}
\it
``We conducted an LLM-based evaluation of requirements utilizing the Llama 2 language model with 70 billion parameters, fine-tuned to complete chat responses...''
\end{quote}

\noindent A more contemporary declaration could similarly state:

\begin{quote}
\it
``We used Claude Opus 4.7 via the Anthropic API to synthesize themes from interview transcripts, with all prompts and conversation logs published as supplementary material.''
\end{quote}

\citeauthor{DBLP:journals/corr/abs-2601-11895}'s DevBench paper illustrates separate disclosure of multiple LLM roles: \emph{GPT-4o} generates the synthetic benchmark instances, nine models (including \emph{Claude 4 Sonnet}, \emph{GPT-4.1}, \emph{DeepSeek-V3}, and \emph{Ministral-3B}) are the evaluation subjects, and \emph{o3-mini} is the LLM judge that scores completions for relevance and helpfulness~\cite{DBLP:journals/corr/abs-2601-11895}.
% GROUNDING DBLP:journals/corr/abs-2601-11895: "Synthetic instances were generated with OpenAI's GPT-4o, chosen for its fluency, reasoning, and code generation capabilities"

\guidelinesubsubsection{Benefits}

Transparency in the use of LLMs helps other researchers understand the context and scope of the study, supporting interpretation and comparison of the results.
Realizing these benefits requires reporting the LLM's exact role and version (see \modelversion).

\guidelinesubsubsection{Challenges}

Declaring LLM usage requires only a brief statement and no additional experiments, making compliance straightforward.
One challenge might be authors' reluctance to disclose LLM usage for valid use cases, because they fear that AI-generated content makes reviewers think that the authors' work is less original.
In fact, there is evidence suggesting that AI disclosure can negatively affect trust in authors~\cite{SCHILKE2025104405}.
% GROUNDING SCHILKE2025104405: "Thirteen experiments consistently demonstrate that actors who disclose their AI usage are trusted less than those who do not."
However, the \emph{ACM Policy on Authorship} requires disclosure of AI used in the research itself, not AI used only to assist with writing~\cite{ACM2026}.
% GROUNDING ACM2026: "ACM no longer requires the disclosure of information regarding the use of AI (as distinct from AI used in the conduct of the research itself"
Our guidelines focus on such use (see \scope), not on proofreading or writing support.

\guidelinesubsubsection{Study Types}

Researchers \must follow this guideline for all study types.
The specific focus of the declaration varies by study type.
For \annotators, \judges, \synthesis, and \subjects, researchers \must declare the specific role assigned to the LLM (e.g., annotator, judge, synthesizer, or simulated participant).
For \llmusage, researchers \must clarify which LLM(s) the observed participants used and under which conditions.
For \newtools, researchers \must declare the LLM's role within the tool architecture and its contribution to the tool's functionality.
For \benchmarkingtasks, researchers \must declare which LLMs were benchmarked and for which tasks.

\guidelinesubsubsection{Advice for Reviewers}

The most common problem with disclosure is incompleteness or vagueness about how the LLM was used. If the paper says ``we used LLM X to help with task Y'' without specifying how, reviewers should request clarification. Such requests are typically minor revisions unless the missing details may reveal methodological problems.

If undisclosed LLM use is suspected, the reviewer should consult their editor or program chair. When the evidence is conclusive, the key question is the degree to which undisclosed use affects the study's contribution, ranging from negligible (e.g., word choice in a single sentence) to severe (e.g., generating the reported data or results).

\guidelinesubsubsection{See Also}
\begin{itemize}[label=$\rightarrow$]
    \item \seemodelversion: The disclosure is incomplete without naming the specific model and version.
    \item \seedesign: When the LLM lives inside a tool or agent, authors must also describe that tool's architecture and prompts.
    \item \seetraces: Session traces show what the LLM did during the study.
\end{itemize}